% =============================================================================
%  Phase 1: Problem Definition & Data Source Identification (project proposal)
%  Limit: 2-3 pages excluding references. Rubric: docs/spec/phase1_details.pdf
%  Build: latexmk -pdf -outdir=build main.tex
% =============================================================================
\documentclass[11pt,letterpaper]{article}
% =============================================================================
%  preamble.tex -- shared styling for all course report deliverables
%  Usage (from reports/phaseN/main.tex):
%      \documentclass[11pt,letterpaper]{article}
%      % =============================================================================
%  preamble.tex -- shared styling for all course report deliverables
%  Usage (from reports/phaseN/main.tex):
%      \documentclass[11pt,letterpaper]{article}
%      % =============================================================================
%  preamble.tex -- shared styling for all course report deliverables
%  Usage (from reports/phaseN/main.tex):
%      \documentclass[11pt,letterpaper]{article}
%      \input{../_template/preamble}
%  Engine: pdfLaTeX + biber (latexmk handles the run order).
%  Do not put content here. Do not use inline colour values in documents;
%  use the named palette below.
% =============================================================================

\makeatletter   % internal \@ names are used throughout; closed at end of file


% -----------------------------------------------------------------------------
%  1. PACKAGES: encoding and language
% -----------------------------------------------------------------------------
\usepackage[T1]{fontenc}          % proper hyphenation of accented glyphs, copyable PDF text
\usepackage[utf8]{inputenc}       % default on modern kernels; kept for older installs
\usepackage[english]{babel}       % hyphenation patterns
\usepackage{csquotes}             % language-aware quotes; biblatex expects it with babel


% -----------------------------------------------------------------------------
%  2. COLOURS: the only place colour values are defined
%     Colour is used for headings, rules, table header rows and links. Nothing else.
% -----------------------------------------------------------------------------
\usepackage[table]{xcolor}        % 'table' loads colortbl for \rowcolor

\definecolor{brandNavy}{HTML}{1F3A5F}    % primary: headings, title, rules
\definecolor{brandAccent}{HTML}{23707D}  % accent: links, title kicker (muted teal)
\definecolor{inkText}{HTML}{1A1D21}      % body text (near-black, softer than pure black)
\definecolor{mutedText}{HTML}{5B6570}    % header/footer, metadata line
\definecolor{ruleGray}{HTML}{B8BEC6}     % light rules (running header)
\definecolor{tableHeadBg}{HTML}{E6ECF2}  % table header row tint (navy at ~10%)


% -----------------------------------------------------------------------------
%  3. FONTS
%     Body:     XCharter, a Type 1 extension of Bitstream Charter. Sturdy, open letterforms
%               that stay legible at 11pt on screen and in print, and a moderate set width.
%     Math:     newtxmath with the xcharter option, so math matches the body.
%     Headings: Source Sans Pro, a humanist sans that pairs with Charter.
%     Mono:     Inconsolata (zi4) for file names and code, scaled to the body x-height.
% -----------------------------------------------------------------------------
\usepackage[osf=false]{XCharter}
\usepackage[xcharter,vvarbb]{newtxmath}
\IfFileExists{sourcesans.sty}            % package renamed in 2022; fall back on older installs
  {\usepackage[semibold]{sourcesans}}    % sets \sffamily only; body stays serif
  {\usepackage[semibold]{sourcesanspro}}
\usepackage[scaled=0.95,varqu]{zi4}      % sets \ttfamily only

\usepackage[final]{microtype}            % protrusion + font expansion: fewer bad breaks, denser text


% -----------------------------------------------------------------------------
%  4. LAYOUT
% -----------------------------------------------------------------------------
\usepackage[letterpaper,margin=1in,headheight=14pt,headsep=12pt,footskip=28pt]{geometry}

\color{inkText}                          % default text colour
\setlength{\parindent}{1.2em}            % indented paragraphs, no gap: densest readable setting
\setlength{\parskip}{0pt}
\linespread{1.04}                        % slight extra leading; Charter is dark and wide

\clubpenalty=10000                       % no widows or orphans
\widowpenalty=10000
\displaywidowpenalty=10000

\usepackage{enumitem}                    % tight lists
\setlist{itemsep=1pt,parsep=0pt,topsep=3pt,partopsep=0pt}
\setlist[itemize]{leftmargin=1.4em}
\setlist[enumerate]{leftmargin=1.8em}
\setlist[description]{leftmargin=1.2em,font=\sffamily\bfseries\color{brandNavy}}

\usepackage{graphicx}
\graphicspath{{../figures/}}             % figures come from reports/figures/ (built by src/.../viz)

\usepackage[font=small,labelfont=bf,skip=4pt]{caption}
\captionsetup[table]{position=top}       % table captions go above the table

\setlength{\textfloatsep}{10pt plus 2pt minus 2pt}   % tighter space around floats
\setlength{\floatsep}{8pt plus 2pt minus 2pt}
\setlength{\intextsep}{8pt plus 2pt minus 2pt}


% -----------------------------------------------------------------------------
%  5. HEADINGS
%     Section:       sans semibold, large, navy, numbered
%     Subsection:    sans semibold, body size, navy
%     Subsubsection: sans semibold, body size, near-black
%     Paragraph:     run-in sans semibold
% -----------------------------------------------------------------------------
\usepackage[explicit]{titlesec}

\titleformat{\section}
  {\sffamily\large\bfseries\color{brandNavy}}{\thesection}{0.6em}{#1}
\titleformat{\subsection}
  {\sffamily\normalsize\bfseries\color{brandNavy}}{\thesubsection}{0.6em}{#1}
\titleformat{\subsubsection}
  {\sffamily\normalsize\bfseries\color{inkText}}{\thesubsubsection}{0.6em}{#1}
\titleformat{\paragraph}[runin]
  {\sffamily\normalsize\bfseries\color{inkText}}{}{0pt}{#1.}

%                         left  before                 after
\titlespacing*{\section}      {0pt}{1.6ex plus .4ex minus .2ex}{0.6ex plus .1ex}
\titlespacing*{\subsection}   {0pt}{1.2ex plus .3ex minus .2ex}{0.4ex plus .1ex}
\titlespacing*{\subsubsection}{0pt}{1.0ex plus .3ex minus .2ex}{0.3ex}
\titlespacing*{\paragraph}    {0pt}{0.8ex plus .2ex minus .1ex}{0.6em}


% -----------------------------------------------------------------------------
%  6. TABLES
%     booktabs rules only (no vertical rules), tinted header row, tabularx for wrapping.
%     Pattern:
%        \tableheadrow \colhead{A} & \colhead{B} \\   <- replaces \toprule + header row
%        \tableheadend                                 <- replaces \midrule under the header
%        body rows ...
%        \bottomrule
%     The header rules have zero separation so the tint meets them with no white gap;
%     \colhead adds symmetric padding instead. Body rows keep standard booktabs spacing.
% -----------------------------------------------------------------------------
\usepackage{array}
\usepackage{booktabs}
\usepackage{tabularx}
\usepackage{ragged2e}

\renewcommand{\arraystretch}{1.15}       % a little air between wrapped body rows
\arrayrulecolor{brandNavy}
\newlength{\cellpad}\setlength{\cellpad}{2.5pt}   % header row padding above and below

% L: ragged-right X column (allows hyphenation, avoids stretched word spacing in narrow cells)
% Y{f}: same, with relative width f. In one table the f values must sum to the number of X/Y columns.
\newcolumntype{L}{>{\RaggedRight\arraybackslash}X}
\newcolumntype{Y}[1]{>{\hsize=#1\hsize\RaggedRight\arraybackslash}X}

\newcommand{\tableheadrow}{\specialrule{\heavyrulewidth}{\abovetopsep}{0pt}\rowcolor{tableHeadBg}}
\newcommand{\tableheadend}{\specialrule{\lightrulewidth}{0pt}{\belowrulesep}}
\newcommand{\colhead}[1]{%
  \rule{0pt}{\dimexpr0.7\baselineskip+\cellpad\relax}%               top padding strut
  {\sffamily\bfseries\color{brandNavy}#1}%
  \rule[\dimexpr-0.3\baselineskip-\cellpad\relax]{0pt}{\dimexpr0.3\baselineskip+\cellpad\relax}}% bottom


% -----------------------------------------------------------------------------
%  7. HEADER / FOOTER
%     Header: short project title (left), phase label (right), thin gray rule.
%     Footer: page number (centre). First page: footer only (title block replaces header).
% -----------------------------------------------------------------------------
\usepackage{fancyhdr}
\pagestyle{fancy}
\fancyhf{}
\fancyhead[L]{\sffamily\small\color{mutedText}\@shorttitle}
\fancyhead[R]{\sffamily\small\color{mutedText}\@phaselabel}
\fancyfoot[C]{\sffamily\small\color{mutedText}\thepage}
\renewcommand{\headrulewidth}{0.4pt}
\renewcommand{\headrule}{{\color{ruleGray}%
  \hrule height\headrulewidth width\headwidth \vskip-\headrulewidth}}

\fancypagestyle{plain}{%
  \fancyhf{}%
  \fancyfoot[C]{\sffamily\small\color{mutedText}\thepage}%
  \renewcommand{\headrulewidth}{0pt}%
}


% -----------------------------------------------------------------------------
%  8. BIBLIOGRAPHY: biblatex + biber, numeric-comp
%     Numeric keeps in-text citations short (matters under a page limit); '-comp'
%     compresses [1,2,3] to [1-3]. Ordered by first citation.
% -----------------------------------------------------------------------------
\usepackage[backend=biber,style=numeric-comp,sorting=none,maxbibnames=6,
            giveninits=true,url=true,doi=true,isbn=false,
            date=iso,urldate=iso,seconds=true]{biblatex}   % ISO dates: unambiguous across locales
\renewcommand*{\bibfont}{\small}
\setlength{\bibitemsep}{2pt}


% -----------------------------------------------------------------------------
%  9. METADATA MACROS
%     Set in each document before \begin{document}:
%       \projecttitle[Short header title]{Full title}
%       \phaselabel{...}  \authors{...}  \course{...}  \reportdate{...}
% -----------------------------------------------------------------------------
\newcommand{\@projecttitle}{[PROJECT TITLE]}
\newcommand{\@shorttitle}{[SHORT TITLE]}
\newcommand{\@phaselabel}{[PHASE LABEL]}
\newcommand{\@authors}{[AUTHORS]}
\newcommand{\@course}{[COURSE]}
\newcommand{\@reportdate}{[DATE]}

\newcommand{\projecttitle}[2][]{%
  \renewcommand{\@projecttitle}{#2}%
  \if\relax\detokenize{#1}\relax\renewcommand{\@shorttitle}{#2}%
  \else\renewcommand{\@shorttitle}{#1}\fi}
\newcommand{\phaselabel}[1]{\renewcommand{\@phaselabel}{#1}}
\newcommand{\authors}[1]{\renewcommand{\@authors}{#1}}
\newcommand{\course}[1]{\renewcommand{\@course}{#1}}
\newcommand{\reportdate}[1]{\renewcommand{\@reportdate}{#1}}

% Compact title block (not a title page). Call right after \begin{document}.
\newcommand{\maketitleblock}{%
  \thispagestyle{plain}%
  \begingroup\raggedright
  {\sffamily\small\bfseries\color{brandAccent}\textls[60]{\MakeUppercase{\@phaselabel}}\par}
  \vspace{3pt}
  {\sffamily\LARGE\bfseries\color{brandNavy}\@projecttitle\par}
  \vspace{6pt}
  {\normalsize\@authors\par}
  \vspace{2pt}
  {\small\color{mutedText}\@course\ \textbar\ \@reportdate\par}
  \vspace{2pt}
  {\color{brandNavy}\rule{\linewidth}{0.8pt}\par}
  \endgroup
  \vspace{4pt}%
}

% Inline code / file names
\newcommand{\code}[1]{\texttt{#1}}


% -----------------------------------------------------------------------------
%  10. HYPERLINKS (load last)
% -----------------------------------------------------------------------------
\usepackage[colorlinks=true,linkcolor=brandNavy,citecolor=brandAccent,
            urlcolor=brandAccent,filecolor=brandAccent,bookmarksnumbered=true,
            pdfstartview=FitH]{hyperref}
\urlstyle{same}                          % URLs in body font: narrower than monospace

\AtBeginDocument{\hypersetup{pdftitle={\@projecttitle},pdfauthor={\@authors},
                             pdfsubject={\@phaselabel}}}
\makeatother

%  Engine: pdfLaTeX + biber (latexmk handles the run order).
%  Do not put content here. Do not use inline colour values in documents;
%  use the named palette below.
% =============================================================================

\makeatletter   % internal \@ names are used throughout; closed at end of file


% -----------------------------------------------------------------------------
%  1. PACKAGES: encoding and language
% -----------------------------------------------------------------------------
\usepackage[T1]{fontenc}          % proper hyphenation of accented glyphs, copyable PDF text
\usepackage[utf8]{inputenc}       % default on modern kernels; kept for older installs
\usepackage[english]{babel}       % hyphenation patterns
\usepackage{csquotes}             % language-aware quotes; biblatex expects it with babel


% -----------------------------------------------------------------------------
%  2. COLOURS: the only place colour values are defined
%     Colour is used for headings, rules, table header rows and links. Nothing else.
% -----------------------------------------------------------------------------
\usepackage[table]{xcolor}        % 'table' loads colortbl for \rowcolor

\definecolor{brandNavy}{HTML}{1F3A5F}    % primary: headings, title, rules
\definecolor{brandAccent}{HTML}{23707D}  % accent: links, title kicker (muted teal)
\definecolor{inkText}{HTML}{1A1D21}      % body text (near-black, softer than pure black)
\definecolor{mutedText}{HTML}{5B6570}    % header/footer, metadata line
\definecolor{ruleGray}{HTML}{B8BEC6}     % light rules (running header)
\definecolor{tableHeadBg}{HTML}{E6ECF2}  % table header row tint (navy at ~10%)


% -----------------------------------------------------------------------------
%  3. FONTS
%     Body:     XCharter, a Type 1 extension of Bitstream Charter. Sturdy, open letterforms
%               that stay legible at 11pt on screen and in print, and a moderate set width.
%     Math:     newtxmath with the xcharter option, so math matches the body.
%     Headings: Source Sans Pro, a humanist sans that pairs with Charter.
%     Mono:     Inconsolata (zi4) for file names and code, scaled to the body x-height.
% -----------------------------------------------------------------------------
\usepackage[osf=false]{XCharter}
\usepackage[xcharter,vvarbb]{newtxmath}
\IfFileExists{sourcesans.sty}            % package renamed in 2022; fall back on older installs
  {\usepackage[semibold]{sourcesans}}    % sets \sffamily only; body stays serif
  {\usepackage[semibold]{sourcesanspro}}
\usepackage[scaled=0.95,varqu]{zi4}      % sets \ttfamily only

\usepackage[final]{microtype}            % protrusion + font expansion: fewer bad breaks, denser text


% -----------------------------------------------------------------------------
%  4. LAYOUT
% -----------------------------------------------------------------------------
\usepackage[letterpaper,margin=1in,headheight=14pt,headsep=12pt,footskip=28pt]{geometry}

\color{inkText}                          % default text colour
\setlength{\parindent}{1.2em}            % indented paragraphs, no gap: densest readable setting
\setlength{\parskip}{0pt}
\linespread{1.04}                        % slight extra leading; Charter is dark and wide

\clubpenalty=10000                       % no widows or orphans
\widowpenalty=10000
\displaywidowpenalty=10000

\usepackage{enumitem}                    % tight lists
\setlist{itemsep=1pt,parsep=0pt,topsep=3pt,partopsep=0pt}
\setlist[itemize]{leftmargin=1.4em}
\setlist[enumerate]{leftmargin=1.8em}
\setlist[description]{leftmargin=1.2em,font=\sffamily\bfseries\color{brandNavy}}

\usepackage{graphicx}
\graphicspath{{../figures/}}             % figures come from reports/figures/ (built by src/.../viz)

\usepackage[font=small,labelfont=bf,skip=4pt]{caption}
\captionsetup[table]{position=top}       % table captions go above the table

\setlength{\textfloatsep}{10pt plus 2pt minus 2pt}   % tighter space around floats
\setlength{\floatsep}{8pt plus 2pt minus 2pt}
\setlength{\intextsep}{8pt plus 2pt minus 2pt}


% -----------------------------------------------------------------------------
%  5. HEADINGS
%     Section:       sans semibold, large, navy, numbered
%     Subsection:    sans semibold, body size, navy
%     Subsubsection: sans semibold, body size, near-black
%     Paragraph:     run-in sans semibold
% -----------------------------------------------------------------------------
\usepackage[explicit]{titlesec}

\titleformat{\section}
  {\sffamily\large\bfseries\color{brandNavy}}{\thesection}{0.6em}{#1}
\titleformat{\subsection}
  {\sffamily\normalsize\bfseries\color{brandNavy}}{\thesubsection}{0.6em}{#1}
\titleformat{\subsubsection}
  {\sffamily\normalsize\bfseries\color{inkText}}{\thesubsubsection}{0.6em}{#1}
\titleformat{\paragraph}[runin]
  {\sffamily\normalsize\bfseries\color{inkText}}{}{0pt}{#1.}

%                         left  before                 after
\titlespacing*{\section}      {0pt}{1.6ex plus .4ex minus .2ex}{0.6ex plus .1ex}
\titlespacing*{\subsection}   {0pt}{1.2ex plus .3ex minus .2ex}{0.4ex plus .1ex}
\titlespacing*{\subsubsection}{0pt}{1.0ex plus .3ex minus .2ex}{0.3ex}
\titlespacing*{\paragraph}    {0pt}{0.8ex plus .2ex minus .1ex}{0.6em}


% -----------------------------------------------------------------------------
%  6. TABLES
%     booktabs rules only (no vertical rules), tinted header row, tabularx for wrapping.
%     Pattern:
%        \tableheadrow \colhead{A} & \colhead{B} \\   <- replaces \toprule + header row
%        \tableheadend                                 <- replaces \midrule under the header
%        body rows ...
%        \bottomrule
%     The header rules have zero separation so the tint meets them with no white gap;
%     \colhead adds symmetric padding instead. Body rows keep standard booktabs spacing.
% -----------------------------------------------------------------------------
\usepackage{array}
\usepackage{booktabs}
\usepackage{tabularx}
\usepackage{ragged2e}

\renewcommand{\arraystretch}{1.15}       % a little air between wrapped body rows
\arrayrulecolor{brandNavy}
\newlength{\cellpad}\setlength{\cellpad}{2.5pt}   % header row padding above and below

% L: ragged-right X column (allows hyphenation, avoids stretched word spacing in narrow cells)
% Y{f}: same, with relative width f. In one table the f values must sum to the number of X/Y columns.
\newcolumntype{L}{>{\RaggedRight\arraybackslash}X}
\newcolumntype{Y}[1]{>{\hsize=#1\hsize\RaggedRight\arraybackslash}X}

\newcommand{\tableheadrow}{\specialrule{\heavyrulewidth}{\abovetopsep}{0pt}\rowcolor{tableHeadBg}}
\newcommand{\tableheadend}{\specialrule{\lightrulewidth}{0pt}{\belowrulesep}}
\newcommand{\colhead}[1]{%
  \rule{0pt}{\dimexpr0.7\baselineskip+\cellpad\relax}%               top padding strut
  {\sffamily\bfseries\color{brandNavy}#1}%
  \rule[\dimexpr-0.3\baselineskip-\cellpad\relax]{0pt}{\dimexpr0.3\baselineskip+\cellpad\relax}}% bottom


% -----------------------------------------------------------------------------
%  7. HEADER / FOOTER
%     Header: short project title (left), phase label (right), thin gray rule.
%     Footer: page number (centre). First page: footer only (title block replaces header).
% -----------------------------------------------------------------------------
\usepackage{fancyhdr}
\pagestyle{fancy}
\fancyhf{}
\fancyhead[L]{\sffamily\small\color{mutedText}\@shorttitle}
\fancyhead[R]{\sffamily\small\color{mutedText}\@phaselabel}
\fancyfoot[C]{\sffamily\small\color{mutedText}\thepage}
\renewcommand{\headrulewidth}{0.4pt}
\renewcommand{\headrule}{{\color{ruleGray}%
  \hrule height\headrulewidth width\headwidth \vskip-\headrulewidth}}

\fancypagestyle{plain}{%
  \fancyhf{}%
  \fancyfoot[C]{\sffamily\small\color{mutedText}\thepage}%
  \renewcommand{\headrulewidth}{0pt}%
}


% -----------------------------------------------------------------------------
%  8. BIBLIOGRAPHY: biblatex + biber, numeric-comp
%     Numeric keeps in-text citations short (matters under a page limit); '-comp'
%     compresses [1,2,3] to [1-3]. Ordered by first citation.
% -----------------------------------------------------------------------------
\usepackage[backend=biber,style=numeric-comp,sorting=none,maxbibnames=6,
            giveninits=true,url=true,doi=true,isbn=false,
            date=iso,urldate=iso,seconds=true]{biblatex}   % ISO dates: unambiguous across locales
\renewcommand*{\bibfont}{\small}
\setlength{\bibitemsep}{2pt}


% -----------------------------------------------------------------------------
%  9. METADATA MACROS
%     Set in each document before \begin{document}:
%       \projecttitle[Short header title]{Full title}
%       \phaselabel{...}  \authors{...}  \course{...}  \reportdate{...}
% -----------------------------------------------------------------------------
\newcommand{\@projecttitle}{[PROJECT TITLE]}
\newcommand{\@shorttitle}{[SHORT TITLE]}
\newcommand{\@phaselabel}{[PHASE LABEL]}
\newcommand{\@authors}{[AUTHORS]}
\newcommand{\@course}{[COURSE]}
\newcommand{\@reportdate}{[DATE]}

\newcommand{\projecttitle}[2][]{%
  \renewcommand{\@projecttitle}{#2}%
  \if\relax\detokenize{#1}\relax\renewcommand{\@shorttitle}{#2}%
  \else\renewcommand{\@shorttitle}{#1}\fi}
\newcommand{\phaselabel}[1]{\renewcommand{\@phaselabel}{#1}}
\newcommand{\authors}[1]{\renewcommand{\@authors}{#1}}
\newcommand{\course}[1]{\renewcommand{\@course}{#1}}
\newcommand{\reportdate}[1]{\renewcommand{\@reportdate}{#1}}

% Compact title block (not a title page). Call right after \begin{document}.
\newcommand{\maketitleblock}{%
  \thispagestyle{plain}%
  \begingroup\raggedright
  {\sffamily\small\bfseries\color{brandAccent}\textls[60]{\MakeUppercase{\@phaselabel}}\par}
  \vspace{3pt}
  {\sffamily\LARGE\bfseries\color{brandNavy}\@projecttitle\par}
  \vspace{6pt}
  {\normalsize\@authors\par}
  \vspace{2pt}
  {\small\color{mutedText}\@course\ \textbar\ \@reportdate\par}
  \vspace{2pt}
  {\color{brandNavy}\rule{\linewidth}{0.8pt}\par}
  \endgroup
  \vspace{4pt}%
}

% Inline code / file names
\newcommand{\code}[1]{\texttt{#1}}


% -----------------------------------------------------------------------------
%  10. HYPERLINKS (load last)
% -----------------------------------------------------------------------------
\usepackage[colorlinks=true,linkcolor=brandNavy,citecolor=brandAccent,
            urlcolor=brandAccent,filecolor=brandAccent,bookmarksnumbered=true,
            pdfstartview=FitH]{hyperref}
\urlstyle{same}                          % URLs in body font: narrower than monospace

\AtBeginDocument{\hypersetup{pdftitle={\@projecttitle},pdfauthor={\@authors},
                             pdfsubject={\@phaselabel}}}
\makeatother

%  Engine: pdfLaTeX + biber (latexmk handles the run order).
%  Do not put content here. Do not use inline colour values in documents;
%  use the named palette below.
% =============================================================================

\makeatletter   % internal \@ names are used throughout; closed at end of file


% -----------------------------------------------------------------------------
%  1. PACKAGES: encoding and language
% -----------------------------------------------------------------------------
\usepackage[T1]{fontenc}          % proper hyphenation of accented glyphs, copyable PDF text
\usepackage[utf8]{inputenc}       % default on modern kernels; kept for older installs
\usepackage[english]{babel}       % hyphenation patterns
\usepackage{csquotes}             % language-aware quotes; biblatex expects it with babel


% -----------------------------------------------------------------------------
%  2. COLOURS: the only place colour values are defined
%     Colour is used for headings, rules, table header rows and links. Nothing else.
% -----------------------------------------------------------------------------
\usepackage[table]{xcolor}        % 'table' loads colortbl for \rowcolor

\definecolor{brandNavy}{HTML}{1F3A5F}    % primary: headings, title, rules
\definecolor{brandAccent}{HTML}{23707D}  % accent: links, title kicker (muted teal)
\definecolor{inkText}{HTML}{1A1D21}      % body text (near-black, softer than pure black)
\definecolor{mutedText}{HTML}{5B6570}    % header/footer, metadata line
\definecolor{ruleGray}{HTML}{B8BEC6}     % light rules (running header)
\definecolor{tableHeadBg}{HTML}{E6ECF2}  % table header row tint (navy at ~10%)


% -----------------------------------------------------------------------------
%  3. FONTS
%     Body:     XCharter, a Type 1 extension of Bitstream Charter. Sturdy, open letterforms
%               that stay legible at 11pt on screen and in print, and a moderate set width.
%     Math:     newtxmath with the xcharter option, so math matches the body.
%     Headings: Source Sans Pro, a humanist sans that pairs with Charter.
%     Mono:     Inconsolata (zi4) for file names and code, scaled to the body x-height.
% -----------------------------------------------------------------------------
\usepackage[osf=false]{XCharter}
\usepackage[xcharter,vvarbb]{newtxmath}
\IfFileExists{sourcesans.sty}            % package renamed in 2022; fall back on older installs
  {\usepackage[semibold]{sourcesans}}    % sets \sffamily only; body stays serif
  {\usepackage[semibold]{sourcesanspro}}
\usepackage[scaled=0.95,varqu]{zi4}      % sets \ttfamily only

\usepackage[final]{microtype}            % protrusion + font expansion: fewer bad breaks, denser text


% -----------------------------------------------------------------------------
%  4. LAYOUT
% -----------------------------------------------------------------------------
\usepackage[letterpaper,margin=1in,headheight=14pt,headsep=12pt,footskip=28pt]{geometry}

\color{inkText}                          % default text colour
\setlength{\parindent}{1.2em}            % indented paragraphs, no gap: densest readable setting
\setlength{\parskip}{0pt}
\linespread{1.04}                        % slight extra leading; Charter is dark and wide

\clubpenalty=10000                       % no widows or orphans
\widowpenalty=10000
\displaywidowpenalty=10000

\usepackage{enumitem}                    % tight lists
\setlist{itemsep=1pt,parsep=0pt,topsep=3pt,partopsep=0pt}
\setlist[itemize]{leftmargin=1.4em}
\setlist[enumerate]{leftmargin=1.8em}
\setlist[description]{leftmargin=1.2em,font=\sffamily\bfseries\color{brandNavy}}

\usepackage{graphicx}
\graphicspath{{../figures/}}             % figures come from reports/figures/ (built by src/.../viz)

\usepackage[font=small,labelfont=bf,skip=4pt]{caption}
\captionsetup[table]{position=top}       % table captions go above the table

\setlength{\textfloatsep}{10pt plus 2pt minus 2pt}   % tighter space around floats
\setlength{\floatsep}{8pt plus 2pt minus 2pt}
\setlength{\intextsep}{8pt plus 2pt minus 2pt}


% -----------------------------------------------------------------------------
%  5. HEADINGS
%     Section:       sans semibold, large, navy, numbered
%     Subsection:    sans semibold, body size, navy
%     Subsubsection: sans semibold, body size, near-black
%     Paragraph:     run-in sans semibold
% -----------------------------------------------------------------------------
\usepackage[explicit]{titlesec}

\titleformat{\section}
  {\sffamily\large\bfseries\color{brandNavy}}{\thesection}{0.6em}{#1}
\titleformat{\subsection}
  {\sffamily\normalsize\bfseries\color{brandNavy}}{\thesubsection}{0.6em}{#1}
\titleformat{\subsubsection}
  {\sffamily\normalsize\bfseries\color{inkText}}{\thesubsubsection}{0.6em}{#1}
\titleformat{\paragraph}[runin]
  {\sffamily\normalsize\bfseries\color{inkText}}{}{0pt}{#1.}

%                         left  before                 after
\titlespacing*{\section}      {0pt}{1.6ex plus .4ex minus .2ex}{0.6ex plus .1ex}
\titlespacing*{\subsection}   {0pt}{1.2ex plus .3ex minus .2ex}{0.4ex plus .1ex}
\titlespacing*{\subsubsection}{0pt}{1.0ex plus .3ex minus .2ex}{0.3ex}
\titlespacing*{\paragraph}    {0pt}{0.8ex plus .2ex minus .1ex}{0.6em}


% -----------------------------------------------------------------------------
%  6. TABLES
%     booktabs rules only (no vertical rules), tinted header row, tabularx for wrapping.
%     Pattern:
%        \tableheadrow \colhead{A} & \colhead{B} \\   <- replaces \toprule + header row
%        \tableheadend                                 <- replaces \midrule under the header
%        body rows ...
%        \bottomrule
%     The header rules have zero separation so the tint meets them with no white gap;
%     \colhead adds symmetric padding instead. Body rows keep standard booktabs spacing.
% -----------------------------------------------------------------------------
\usepackage{array}
\usepackage{booktabs}
\usepackage{tabularx}
\usepackage{ragged2e}

\renewcommand{\arraystretch}{1.15}       % a little air between wrapped body rows
\arrayrulecolor{brandNavy}
\newlength{\cellpad}\setlength{\cellpad}{2.5pt}   % header row padding above and below

% L: ragged-right X column (allows hyphenation, avoids stretched word spacing in narrow cells)
% Y{f}: same, with relative width f. In one table the f values must sum to the number of X/Y columns.
\newcolumntype{L}{>{\RaggedRight\arraybackslash}X}
\newcolumntype{Y}[1]{>{\hsize=#1\hsize\RaggedRight\arraybackslash}X}

\newcommand{\tableheadrow}{\specialrule{\heavyrulewidth}{\abovetopsep}{0pt}\rowcolor{tableHeadBg}}
\newcommand{\tableheadend}{\specialrule{\lightrulewidth}{0pt}{\belowrulesep}}
\newcommand{\colhead}[1]{%
  \rule{0pt}{\dimexpr0.7\baselineskip+\cellpad\relax}%               top padding strut
  {\sffamily\bfseries\color{brandNavy}#1}%
  \rule[\dimexpr-0.3\baselineskip-\cellpad\relax]{0pt}{\dimexpr0.3\baselineskip+\cellpad\relax}}% bottom


% -----------------------------------------------------------------------------
%  7. HEADER / FOOTER
%     Header: short project title (left), phase label (right), thin gray rule.
%     Footer: page number (centre). First page: footer only (title block replaces header).
% -----------------------------------------------------------------------------
\usepackage{fancyhdr}
\pagestyle{fancy}
\fancyhf{}
\fancyhead[L]{\sffamily\small\color{mutedText}\@shorttitle}
\fancyhead[R]{\sffamily\small\color{mutedText}\@phaselabel}
\fancyfoot[C]{\sffamily\small\color{mutedText}\thepage}
\renewcommand{\headrulewidth}{0.4pt}
\renewcommand{\headrule}{{\color{ruleGray}%
  \hrule height\headrulewidth width\headwidth \vskip-\headrulewidth}}

\fancypagestyle{plain}{%
  \fancyhf{}%
  \fancyfoot[C]{\sffamily\small\color{mutedText}\thepage}%
  \renewcommand{\headrulewidth}{0pt}%
}


% -----------------------------------------------------------------------------
%  8. BIBLIOGRAPHY: biblatex + biber, numeric-comp
%     Numeric keeps in-text citations short (matters under a page limit); '-comp'
%     compresses [1,2,3] to [1-3]. Ordered by first citation.
% -----------------------------------------------------------------------------
\usepackage[backend=biber,style=numeric-comp,sorting=none,maxbibnames=6,
            giveninits=true,url=true,doi=true,isbn=false,
            date=iso,urldate=iso,seconds=true]{biblatex}   % ISO dates: unambiguous across locales
\renewcommand*{\bibfont}{\small}
\setlength{\bibitemsep}{2pt}


% -----------------------------------------------------------------------------
%  9. METADATA MACROS
%     Set in each document before \begin{document}:
%       \projecttitle[Short header title]{Full title}
%       \phaselabel{...}  \authors{...}  \course{...}  \reportdate{...}
% -----------------------------------------------------------------------------
\newcommand{\@projecttitle}{[PROJECT TITLE]}
\newcommand{\@shorttitle}{[SHORT TITLE]}
\newcommand{\@phaselabel}{[PHASE LABEL]}
\newcommand{\@authors}{[AUTHORS]}
\newcommand{\@course}{[COURSE]}
\newcommand{\@reportdate}{[DATE]}

\newcommand{\projecttitle}[2][]{%
  \renewcommand{\@projecttitle}{#2}%
  \if\relax\detokenize{#1}\relax\renewcommand{\@shorttitle}{#2}%
  \else\renewcommand{\@shorttitle}{#1}\fi}
\newcommand{\phaselabel}[1]{\renewcommand{\@phaselabel}{#1}}
\newcommand{\authors}[1]{\renewcommand{\@authors}{#1}}
\newcommand{\course}[1]{\renewcommand{\@course}{#1}}
\newcommand{\reportdate}[1]{\renewcommand{\@reportdate}{#1}}

% Compact title block (not a title page). Call right after \begin{document}.
\newcommand{\maketitleblock}{%
  \thispagestyle{plain}%
  \begingroup\raggedright
  {\sffamily\small\bfseries\color{brandAccent}\textls[60]{\MakeUppercase{\@phaselabel}}\par}
  \vspace{3pt}
  {\sffamily\LARGE\bfseries\color{brandNavy}\@projecttitle\par}
  \vspace{6pt}
  {\normalsize\@authors\par}
  \vspace{2pt}
  {\small\color{mutedText}\@course\ \textbar\ \@reportdate\par}
  \vspace{2pt}
  {\color{brandNavy}\rule{\linewidth}{0.8pt}\par}
  \endgroup
  \vspace{4pt}%
}

% Inline code / file names
\newcommand{\code}[1]{\texttt{#1}}


% -----------------------------------------------------------------------------
%  10. HYPERLINKS (load last)
% -----------------------------------------------------------------------------
\usepackage[colorlinks=true,linkcolor=brandNavy,citecolor=brandAccent,
            urlcolor=brandAccent,filecolor=brandAccent,bookmarksnumbered=true,
            pdfstartview=FitH]{hyperref}
\urlstyle{same}                          % URLs in body font: narrower than monospace

\AtBeginDocument{\hypersetup{pdftitle={\@projecttitle},pdfauthor={\@authors},
                             pdfsubject={\@phaselabel}}}
\makeatother

\addbibresource{references.bib}

% ---- Metadata ---------------------------------------------------------------
\projecttitle[Power Consumption Forecasting]{Power Consumption Forecasting Using Smart Meter Data}
\phaselabel{Phase 1 \textperiodcentered{} Project Proposal}
\authors{[AUTHOR 1], [AUTHOR 2], [AUTHOR 3], [AUTHOR 4]}
\course{Data Analytics, Master of Computer Science, University of New Brunswick}
\reportdate{[DATE]}

\begin{document}
\maketitleblock

% Submission checklist (rubric, "Deliverable & Submission Requirements"):
%   - Clear project title and the names of ALL team members (title block).
%   - Headings match the five proposal components (sections below).
%   - Cite external datasets, APIs, reports used to justify the project.
%   - Concise, professional writing.
%   - Clearly distinguish CONFIRMED data access from PLANNED / potential sources.
%   - 2-3 pages excluding references/appendices. One proposal per team.


% =============================================================================
\section{Problem Statement}
\label{sec:problem}
% Rubric requires:
%   - The business, organizational, or research problem being investigated.
%   - Context: why the problem matters.
%   - The population, process, customers, organization, or system studied.
%   - The problem stated in a specific, MEASURABLE way, not a broad topic.
% Quality standard (strong): specific, meaningful, measurable problem.


% =============================================================================
\section{Project Objectives \& Analytics Goal}
\label{sec:objectives}
% Rubric requires:
%   - 2-4 specific objectives (list below).
%   - The primary analytics task (prediction, classification, clustering,
%     visualization, anomaly detection, or other).
%   - The target/outcome to be predicted or explained, if applicable.
%   - What a successful project would produce or enable.
% Quality standard (strong): clear task and intended outcome.

\begin{enumerate}
  \item {[Objective 1]}
  \item {[Objective 2]}
  % \item {[Objective 3]}
  % \item {[Objective 4]}
\end{enumerate}


% =============================================================================
\section{Data Source Identification}
\label{sec:sources}
% Rubric requires, for each main data source:
%   - What information it contains and why it is relevant to the problem.
%   - Its type: internal, external, API, sensor/device, survey, transactional,
%     public/open data, or other.
%   - Expected format and approximate size or time period, when known.
% Status column: mark each source Confirmed or Planned (submission requirement).
% Quality standard (strong): relevant, credible, feasible sources.

\begin{table}[htbp]
  \centering
  \caption{[Caption]}
  \label{tab:sources}
  % Y{f} widths must sum to 5 (the number of columns).
  \begin{tabularx}{\linewidth}{Y{0.95} Y{0.70} Y{1.45} Y{1.20} Y{0.70}}
    \tableheadrow
    \colhead{Source} & \colhead{Type} & \colhead{What it provides}
      & \colhead{Access / feasibility} & \colhead{Status} \\
    \tableheadend
    % Brace a cell that starts with [ right after a rule: rules take an optional [width].
    {[Source]} & [Type] & [What it provides] & [Access / feasibility] & [Status] \\
    \bottomrule
  \end{tabularx}
\end{table}


% =============================================================================
\section{Data Access, Quality \& Feasibility}
\label{sec:feasibility}
% Rubric requires:
%   - How the data will be obtained; whether access is already available.
%   - Obvious issues: missing values, inconsistent records, limited
%     observations, bias, data integration challenges.
%   - Privacy, security, consent, licensing, or ethical considerations.
%   - Why the sources are sufficient, or what additional data may be needed.
% Quality standard (strong): access, quality, ethics, and constraints
% considered; no major data risk overlooked.


% =============================================================================
\section{Preliminary Project Scope}
\label{sec:scope}
% Rubric requires:
%   - What the project WILL and WILL NOT attempt.
%   - Expected users or stakeholders of the final analysis.
%   - Expected final outputs (dashboard, analytical report, predictive model,
%     decision-support tool, ...).


% =============================================================================
% References do not count toward the page limit.
\printbibliography[title={References}]

\end{document}
